\begin{trapbox}
\begin{description}[style=nextline, leftmargin=2.8em, labelsep=0.5em, itemsep=0.45em]
\item[\textbf{\textcolor{CTrap}{易错点一：把焦点当成顶点}}]
$F_1(-c,0),F_2(c,0)$ 是\textbf{焦点}；$A_1(-a,0),A_2(a,0)$ 是长轴顶点。“两焦点到切线的距离乘积等于 $b^2$”不能擅自换成“两顶点”。

\item[\textbf{\textcolor{CTrap}{易错点二：距离乘积反推相切，丢了绝对值}}]
对\textbf{任意直线} $y=kx+m$，只有
\[
 d_1d_2=\frac{|m^2-c^2k^2|}{1+k^2}.
\]
单令 $d_1d_2=b^2$ 会出现额外分支
\[
 m^2=c^2k^2-b^2(1+k^2),
\]
它可能是割线，并非切线。\textbf{补足“两个焦点同侧”即 $m^2>c^2k^2$ 才能去掉绝对值。}反例：$x^2/9+y^2/4=1$ 与 $y=2x$，乘积为 $4$ 却经过椭圆中心。

\item[\textbf{\textcolor{CTrap}{易错点三：公式里的 $a^2,b^2$ 对调或写成减号}}]
对 $x^2/a^2+y^2/b^2=1$ 和 $y=kx+m$，一定是
\[
 \boxed{m^2=a^2k^2+b^2}.
\]
口诀：\textbf{斜率平方乘 $x$ 分母，常数项取 $y$ 分母。}若写成 $x^2/A^2+y^2/B^2=1$，即便 $A<B$，代数公式仍为 $m^2=A^2k^2+B^2$；只是焦点的方向随之改变。

\item[\textbf{\textcolor{CTrap}{易错点四：误解截距太大／太小}}]
固定 $k$，$|m|<\sqrt{a^2k^2+b^2}$ 是\textbf{割线}；$|m|=\sqrt{a^2k^2+b^2}$ 是\textbf{切线}；$|m|>\sqrt{a^2k^2+b^2}$ 是\textbf{相离}。特别地，$|m|<b$ 时直线经过椭圆内部的 $(0,m)$，不可能相切。原例 $m=2,b=\sqrt5$ 是\textbf{截距绝对值太小}。

\item[\textbf{\textcolor{CTrap}{易错点五：把 $m^2=\cdots$ 当作唯一一条切线}}]
固定 $k$ 时有 $\boxed{m=\pm\sqrt{a^2k^2+b^2}}$，通常是\textbf{两条平行切线}；若题目另有限定（如切点位于上半平面），还要挑选符号。

\item[\textbf{\textcolor{CTrap}{易错点六：遗漏竖直切线}}]
$y=kx+m$ 只包括斜率存在的直线。两条竖直切线 $x=\pm a$ 不在本公式的取值范围内，不能用“$k$ 等于无穷大”直接代公式。

\item[\textbf{\textcolor{CTrap}{易错点七：误将点到直线距离当作 $|m|$}}]
对于 $kx-y+m=0$，点 $(x_0,y_0)$ 的距离是
\[
 \frac{|kx_0-y_0+m|}{\sqrt{k^2+1}}.
\]
尤其两个焦点的距离分子分别为 $|m-kc|$、$|m+kc|$，\textbf{分母必须是 $\sqrt{1+k^2}$}。

\item[\textbf{\textcolor{CTrap}{易错点八：切点公式不核验条件}}]
$P=(-a^2k/m,b^2/m)$ 是\textbf{相切成立时}的切点公式。若直线并未满足 $m^2=a^2k^2+b^2$，不能直接称其为切点；在相切时 $m\ne0$，所以分母合法。
\end{description}
\end{trapbox}
